% Appendix D — Information-parity round-trip validation
% Measured recovery rates from the round-trip validation.

\section{Information-parity round-trip validation}
\label{app:D}

\subsection*{Method}

To quantify how much of the Core IR (C4) survives the projection to the
information-parity condition (C1+), we run a round trip
$\text{C4} \rightarrow \text{C1+} \rightarrow \text{C4}'$ on all 30 retained
confirmatory functions and compare C4 against the re-derived C4$'$. C1+ is
produced by the same projection used in the experiment
(\texttt{src/pipeline/ir\_to\_annotated\_text.py}). C4$'$ is re-derived from the
C1+ \emph{text only}: the AST by re-parsing the C1+ source, the contract
clauses by parsing the structured docstring block (a docstring$\rightarrow$contract
parser written for this purpose, since the production lowering reads the
authoritative reviewed contract rather than the projected text), and the types
from the C1+ annotations. Comparison normalization is specified in
\texttt{docs/roundtrip-comparison-spec.md}; the implementation is
\texttt{src/experiment/roundtrip.py} and the per-function and aggregate outputs
are under \texttt{results/analysis/exploratory/roundtrip/}. This is an
exploratory measurement, not a confirmatory test.

\subsection*{Results}

Table~\ref{tab:roundtrip} reports per-metric recovery over the 30 functions.

\begin{table}[h]
\centering
\small
\caption{C4$\rightarrow$C1+$\rightarrow$C4$'$ recovery over the 30 retained
confirmatory functions (\emph{Exploratory analysis}). Retention is the fraction
of the C4 quantity recovered in C4$'$; see
\texttt{docs/roundtrip-comparison-spec.md} for the normalization of each metric.}
\label{tab:roundtrip}
\begin{tabular}{lrrr}
\toprule
Metric & Mean & Min & Max \\
\midrule
Node-kind retention            & $0.998$ & $0.962$ & $1.000$ \\
Type-annotation retention      & $0.964$ & $0.600$ & $1.000$ \\
Parent--child (edge) retention & $0.987$ & $0.885$ & $1.000$ \\
Cross-reference retention      & $1.000$ & $1.000$ & $1.000$ \\
Contract preconditions         & $1.000$ & $1.000$ & $1.000$ \\
Contract postconditions        & $1.000$ & $1.000$ & $1.000$ \\
Contract invariants            & $1.000$ & $1.000$ & $1.000$ \\
\midrule
Canonical-hash agreement       & \multicolumn{3}{l}{$0/30$ functions (exact AST match)} \\
\bottomrule
\end{tabular}
\end{table}

The propositional content the experiment relies on is recovered essentially in
full: all three contract-clause families and the cross-reference structure among
clauses are recovered at $100\%$, node kinds at $99.8\%$, type annotations at
$96.4\%$, and parent--child structure at $98.7\%$. What is \emph{not} recovered
is byte-exact AST identity: $0/30$ functions round-trip to an identical AST,
because C1+ introduces a structured docstring node and does not preserve the
IR's syntactic child-ordering or node-level metadata.

\subsection*{Interpretation}

These rates bear on the information-parity claim but do not settle the study's
question, and must not be read in a confirmatory tone. High recovery means C1+
is a high-fidelity, information-preserving projection of C4 rather than a lossy
summary: the contract content and cross-references that the reasoning task
depends on survive the projection intact. \emph{If} C4 is nonetheless found to
outperform C1+ in a future adequately-powered comparison, this high recovery
\emph{strengthens}---but does not establish---the reading that representational
structure itself, not differing information content, drives the effect; the
residual non-recovered structure (byte-exact AST identity, $3.6\%$ of type
annotations, $1.3\%$ of edges) remains a candidate confound. Conversely, were
recovery low on the content metrics, the parity claim would be weakened and the
C4-vs-C1+ contrast would be partly confounded by information loss; that is not
what we observe here for the content metrics, but the never-exact AST identity
is stated plainly as the boundary of the parity claim. No causal or confirmatory
conclusion is drawn from this exploratory measurement.
